\chapter{A Concrete Naming Certificate for $\BishopIntervalNestingMetricCompletionUp$}
\label{ch:concrete-instances-bishop-interval-nesting-metric-completion-namecert}
\origin{human}

Following Bishop and Bridges on nested intervals and constructive metric completion, the $\BishopIntervalNestingMetricCompletionUp$ packet records the handoff from displayed interval nesting to a metric-completion-facing real name as finite BEDC data. It sits after the located nested-interval route of \autoref{ch:concrete-instances-located-real-nested-intervals-namecert}, the dyadic arithmetic core of \autoref{ch:concrete-instances-dyadicratcore-namecert}, the finite stream interface of \autoref{ch:concrete-instances-streamname-namecert}, the regular rational readback of \autoref{ch:concrete-instances-regseqrat-namecert}, the metric completion boundary of \autoref{ch:concrete-instances-metriccompletion-namecert}, and the terminal real-name seal of \autoref{ch:concrete-instances-real-namecert}. The packet names only the displayed nesting-to-completion carrier; it does not assert ambient $\RealUp$ completeness, choose an interval intersection point, quotient endpoint names, or select a metric-completion representative.

\begin{definition}[Bishop interval nesting metric-completion carrier]
\label{def:bishop-interval-nesting-metric-completion-carrier}
A $\BishopIntervalNestingMetricCompletionUp$ carrier is a finite $\BHist$ packet
$$
K=(I,D,S,N,W,R,M,E,H,C,P,L)
$$
with the following visible rows:
\begin{enumerate}
\item $I$ is the located interval row, displaying lower and upper endpoint witnesses and the locatedness comparison needed for retained windows.
\item $D$ is the dyadic bisection and midpoint row, read through $\DyadicRatCoreUp$ and carrying the parent-to-child radius-halving ledger.
\item $S$ is the shrinking-diameter budget row recording the finite request-to-radius schedule.
\item $N$ is the nested containment ledger, including the retained child relation and the proof that later windows remain inside earlier windows.
\item $W$ is the $\StreamNameUp$ finite-window row for the retained interval sequence.
\item $R$ is the $\RegSeqRatUp$ readback row exposed after the stream window and dyadic tolerance data have been displayed.
\item $M$ is the $\MetricCompletionUp$ handoff row through which completion-facing consumers read the same regular Cauchy route without choosing a quotient representative.
\item $E$ is the terminal $\RealUp$ seal row, downstream of $W,R,M$.
\item $H$ records componentwise $\hsame$ transport for the interval, bisection, diameter, nesting, stream, readback, completion, and seal rows.
\item $C$ records displayed $\Cont$ replay from a precision request through retained windows and into completion-facing readback.
\item $P$ is the $\Pkg$ provenance over $\BishopIntervalNestingMetricCompletionUp$, $\LocatedRealNestedIntervalsUp$, $\DyadicRatCoreUp$, $\StreamNameUp$, $\RegSeqRatUp$, $\MetricCompletionUp$, $\RealUp$, $\BHist$, $\hsame$, $\Cont$, and $\NameCert$ rows.
\item $L$ is the local $\NameCert_{\BishopIntervalNestingMetricCompletionUp}$ row.
\end{enumerate}
The carrier contains no coordinate for an ambient completeness theorem, a selected interval-intersection point, quotient equality of endpoint names, countable choice, trichotomy, or a detached real seal.
\end{definition}

\begin{theorem}[Bishop interval nesting metric-completion obligations]
\label{thm:bishop-interval-nesting-metric-completion-namecert-obligations}
For every accepted $\BishopIntervalNestingMetricCompletionUp$ carrier $K=(I,D,S,N,W,R,M,E,H,C,P,L)$, the local $\NameCert_{\BishopIntervalNestingMetricCompletionUp}$ surface consists exactly of these finite obligations:
\begin{enumerate}
\item carrier admission for the displayed located interval, dyadic midpoint, shrinking-diameter, containment, stream, readback, completion, real-seal, transport, replay, provenance, and local naming rows;
\item dyadic contraction, requiring each retained child interval to read the same $\DyadicRatCoreUp$ midpoint and radius-halving ledger before the next window is admitted;
\item nested containment, requiring the row $N$ to display that every later retained interval remains inside the earlier window named by the same stream prefix;
\item regular readback, requiring $W$ and $R$ to expose the retained windows as a $\RegSeqRatUp$ route before $M$ or $E$ is read;
\item metric-completion handoff, requiring every $\MetricCompletionUp$ consumer to read $I,D,S,N,W,R$ on the same packet before using $M$;
\item real-seal non-escape, requiring $E$ to be downstream of the displayed stream, regular-readback, and completion rows and excluding all hidden completeness or quotient coordinates.
\end{enumerate}
\end{theorem}
\begin{proof}
Project the carrier of \autoref{def:bishop-interval-nesting-metric-completion-carrier}. This exposes exactly the rows $I,D,S,N,W,R,M,E,H,C,P,L$ and no additional host-level interval intersection, quotient representative, choice functional, trichotomy, or completeness coordinate.

The dyadic contraction obligation is read from $D$ and $S$. The midpoint and radius-halving data are carried as dyadic rows, while the shrinking schedule tells which finite budget is attached to the retained child interval.

The nesting and readback obligations are then rowwise. The containment ledger $N$ records inclusion between retained windows, $W$ names the finite stream windows, and $R$ repacks those windows as regular rational readbacks.

Finally, $M$ and $E$ are downstream rows on the same packet. The transport, replay, provenance, and local naming rows $H,C,P,L$ preserve the route, so every completion-facing or real-seal read is forced through the displayed finite obligations and cannot export an excluded host object.
\end{proof}

\begin{theorem}[Bishop interval nesting metric-completion point seal]
\label{thm:bishop-interval-nesting-metric-completion-point-seal}
Every $\RealUp$- or $\MetricCompletionUp$-facing read of an accepted $\BishopIntervalNestingMetricCompletionUp$ carrier factors through
$$
I,D,S,N,W,R,M,E,H,C,P,L.
$$
In particular, the point seal is the displayed route from located nested intervals through $\StreamNameUp$ windows, $\RegSeqRatUp$ readback, and the $\MetricCompletionUp$ handoff before the terminal $\RealUp$ row is exposed; it is not an ambient intersection theorem or a selected completed point.
\end{theorem}
\begin{proof}
Begin with the accepted carrier and read its located interval row $I$. The dyadic row $D$ supplies midpoint and child-window data, while $S$ supplies the shrinking diameter budget that controls the requested precision.

The retained windows cannot be consumed until the nested containment ledger $N$ is displayed. Once $N$ is present, the stream row $W$ names the finite window sequence and the regular row $R$ provides the $\RegSeqRatUp$ readback seen by downstream real and completion consumers.

The metric-completion handoff $M$ is therefore not a selected quotient representative; it is the same finite regular route repackaged for $\MetricCompletionUp$. The real seal $E$ is downstream of that handoff and remains transported by $H,C,P,L$.

Since the carrier has no field for an ambient completeness theorem, selected intersection point, quotient endpoint equality, countable choice, or detached real seal, every point-seal read is exactly the displayed finite route.
\end{proof}

\closureat{BishopIntervalNestingMetricCompletionUp}{seedStr}

\begin{closurestatus}{\BishopIntervalNestingMetricCompletionUp}
  \constructivestory{A $\BishopIntervalNestingMetricCompletionUp$ packet is generated from located interval rows, dyadic midpoint and bisection rows, shrinking diameter budgets, nested containment ledgers, StreamName windows, RegSeqRat readback, a MetricCompletion handoff, a terminal Real seal, componentwise hsame transport, Cont replay, Pkg provenance, and a local NameCert row.}
  \theoryclosure{\seedClosure}
  \scopeclosed{\autoref{def:bishop-interval-nesting-metric-completion-carrier}, \autoref{thm:bishop-interval-nesting-metric-completion-namecert-obligations}, and \autoref{thm:bishop-interval-nesting-metric-completion-point-seal} seed the finite handoff from Bishop-style interval nesting to MetricCompletion and Real consumers through displayed dyadic, stream, and regular-readback rows.}
  \formalstatus{\unformalizedV}
  \bridgestatus{none}
  \notclaimed{This chapter does not close ambient Real completeness, selected interval intersections, quotient equality of endpoint names, countable choice, trichotomy, Lean verification, or bridge checking.}
  \upgradepath{Scoped closure requires checked carrier admission, dyadic contraction, nested containment, StreamName window exposure, RegSeqRat readback, MetricCompletion handoff, Real seal nonescape, transport, replay, provenance, and local naming for BEDC.Derived.BishopIntervalNestingMetricCompletionUp.}
\end{closurestatus}
